\chapter{Declaración de uso de inteligencia artificial}
\label{apendice:declaracionIA}

Este anexo detalla el uso que he hecho de herramientas de inteligencia artificial generativa durante el trabajo, y complementa la declaración responsable sobre autoría y uso ético de herramientas de inteligencia artificial que la Universidad Complutense de Madrid pide firmar con la entrega. Recoge qué herramientas y modelos usé y en qué periodos, para qué tareas, y qué parte del trabajo se hizo sin ellas.

\section{Herramientas y modelos}
\label{sec:declaracionHerramientas}

\begin{itemize}
    \item \textbf{Claude Code (Anthropic).} Ha sido la herramienta principal desde enero de 2026 hasta la entrega, en septiembre de 2026. Es el asistente de programación de Anthropic para el terminal y el editor, y trabaja directamente sobre el repositorio del proyecto. En cada momento usé la versión de la familia Claude Opus disponible en la herramienta y, desde junio de 2026, también la familia Claude Fable. Con los identificadores de la documentación de Anthropic~\citep{anthropic2026models}, los modelos publicados en ese periodo fueron Claude Opus 4.5 (\texttt{claude-opus-4-5-20251101}, noviembre de 2025, el vigente en enero de 2026), Claude Opus 4.6 (\texttt{claude-opus-4-6}, febrero de 2026), Claude Opus 4.7 (\texttt{claude-opus-4-7}, abril de 2026), Claude Opus 4.8 (\texttt{claude-opus-4-8}, mayo de 2026), Claude Fable 5 (\texttt{claude-fable-5}, junio de 2026), Claude Opus 5 (\texttt{claude-opus-5}, julio de 2026) y Claude Fable 5.1 (\texttt{claude-fable-5-1}, septiembre de 2026). No conservo un registro de qué modelo atendió cada sesión, así que declaro el rango completo.
    \item \textbf{Figma desde Claude Code.} El prototipo navegable de Figma se construyó y se actualizó desde Claude Code a través del servidor MCP oficial de Figma (\textit{Model Context Protocol}), que permite al asistente leer y modificar el fichero de diseño. Las pantallas del segundo prototipo y su actualización tras la auditoría entre diseño e implementación (capítulos~\ref{cap:analisisRequisitos} y~\ref{cap:disenoSistema}) se hicieron así, sobre mis indicaciones y con mi revisión en Figma.
    \item \textbf{Gemini (Google).} Antes de Claude Code, en junio y julio de 2025 y en diciembre de 2025, usé la aplicación Gemini de Google con los modelos que ofrecía entonces: en junio y julio de 2025 la familia Gemini 2.5 (Gemini 2.5 Pro y Gemini 2.5 Flash, publicados como versiones estables ese mismo junio) y en diciembre de 2025 la familia Gemini 3 (Gemini 3 Pro y Gemini 3 Flash, publicados en noviembre y diciembre de 2025) junto con la 2.5~\citep{google2026gemini}. No conservo registro de qué variante respondió a cada consulta.
\end{itemize}

No he usado ChatGPT ni ninguna otra herramienta de inteligencia artificial generativa.

\section{Tareas en las que se usaron}
\label{sec:declaracionTareas}

% Victor: confirmar para que usaste Gemini en 2025 (regla 21); el chat no lo sabe.
Gemini lo usé para consultas puntuales de apoyo mientras definía la idea y la primera versión de la propuesta, la que recoge el Anexo~\ref{apendice:catalogoV1}.

Claude Code intervino en cinco tipos de tarea:

\begin{itemize}
    \item \textbf{Código.} La mayor parte del código del repositorio (backend en Python con el motor y el copiloto, frontend, migraciones SQL y scripts de experimentos) se escribió con su asistencia a partir de mis especificaciones, bloque a bloque. Cada bloque tiene un banco de pruebas permanente que se ejecuta antes de cerrarlo y con cada cambio posterior, y revisé cada cambio antes de incorporarlo al repositorio.
    \item \textbf{Diseño.} Las pantallas de Figma descritas en el punto anterior.
    \item \textbf{Experimentos.} El diseño de los experimentos del capítulo~\ref{cap:evaluacion}, con sus criterios fijados por escrito antes de medir, y los scripts que los ejecutan y producen sus tablas y figuras se prepararon con la herramienta. Las cifras son las que devolvieron esos scripts sobre el sistema real.
    \item \textbf{Memoria.} Los capítulos de esta memoria partieron de borradores redactados con Claude Code a partir de la documentación interna del proyecto y de mis indicaciones sobre qué contar y cómo. Los revisé y edité uno a uno; mis ediciones sobre los primeros borradores sirvieron de referencia de estilo para los siguientes, y el texto final es el que aprobé párrafo a párrafo. Las traducciones al inglés del resumen, la introducción y las conclusiones se hicieron de la misma forma.
    \item \textbf{Documentación interna.} Las especificaciones de cada bloque, el registro de decisiones y las notas de las reuniones.
\end{itemize}

\section{Autoría y verificación}
\label{sec:declaracionAutoria}

Las decisiones de alcance, de arquitectura y de modelado son mías, tomadas en las reuniones con los directores del TFG y a partir de lo que pedían los nutricionistas consultados. Las entrevistas con profesionales y la prueba con usuarios las hice yo. Los datos y cifras de la evaluación están medidos sobre el sistema, no generados por ningún modelo. La bibliografía está verificada en su fuente. Y ninguna parte del sistema desplegado ni de esta memoria se ha incorporado sin mi revisión.
